% !TEX program = xelatex
\documentclass[UTF8,a4paper,11pt]{ctexart}
\usepackage[a4paper,margin=2.35cm]{geometry}
\usepackage{amsmath,amssymb,mathtools}
\usepackage{booktabs,longtable,tabularx,array,multirow}
\usepackage{graphicx,xcolor,hyperref,cleveref}
\usepackage{tikz}
\usetikzlibrary{arrows.meta,positioning,fit,calc,shapes.geometric}
\usepackage{enumitem}
\usepackage{microtype}
\usepackage{siunitx}
\usepackage{algorithm}
\usepackage{algpseudocode}
\hypersetup{colorlinks=true,linkcolor=blue!60!black,citecolor=blue!60!black,urlcolor=blue!60!black,pdftitle={Bridge-RGS}}
\setlist{nosep,leftmargin=2em}
\setlength{\emergencystretch}{3em}
\newcommand{\method}{\textsc{Bridge-RGS}}
\newcommand{\DINO}{\textsc{DINOv3}}
\newcommand{\GS}{3D Gaussian Splatting}
\newcommand{\mIoU}{\ensuremath{\mathrm{mIoU}}}
\newcommand{\PSNR}{\ensuremath{\mathrm{PSNR}}}
\newcommand{\LPIPS}{\ensuremath{\mathrm{LPIPS}}}
\newcommand{\E}{\mathbb{E}}
\newcommand{\R}{\mathbb{R}}
\title{Bridge-RGS：相机鲁棒、结构自适应的桥梁 RGB--语义联合高斯重建}
\author{面向 IC-SHM 2026 Project 2 的可复现实验报告}
\date{\today}
\begin{document}
\maketitle
\begin{abstract}
桥梁巡检场景要求从有限相机轨迹同时恢复可视化质量较高的 RGB 场景和官方五类语义掩码。本文报告 \method{}，一种由“语义场内部共享几何”和“最终多场部署组合”组成的桥梁 Gaussian 重建系统。语义场以 \GS{} 为几何表示，以 ModelScope 发布的高容量 \DINO{} ViT-H+/16 作为部署阶段的冻结预训练视觉编码器，并由可训练语义解码器将其特征转换为类别概率；两套独立 RGB 场负责最终 RGB 质量。研究框架还包含投影可靠性、多视图语义融合、归一化梯度增密和相机残差诊断等可选模块，本文严格区分实际采用路径与历史候选。

在固定官方传感器、50 张 RGB 留出视角和其中 41 张语义标注视角的共同原图协议上，最终三场相机条件联合系统达到 \PSNR{}=30.236593\,dB、SSIM=0.875024、\LPIPS{}=0.259300 和五类 \mIoU{}=95.109016\%。相对本地同协议 Swin-L 参考，PSNR 提高 0.835644\,dB，\mIoU{} 提高 0.790633 个百分点，且三个 RGB 指标均改善。相对基准三场联合系统，PSNR 进一步提高 0.200932\,dB，语义混淆矩阵保持不变。报告也给出与外部进度目录的协议边界、计算成本、失败模式和复现步骤。上述结果支持当前工程交付，不等同于跨场景或同预算的赛事盲测胜出；学术创新应由严格机制消融继续验证。
\end{abstract}
\noindent\textbf{关键词：}桥梁巡检；3D Gaussian Splatting；DINOv3；语义渲染；相机不确定性；多视图融合；可复现评测

\section{引言}
桥梁巡检图像同时包含宽阔背景、细长拉索、桥塔和局部基础结构。对 RGB 重建而言，纹理、曝光和高频边界决定感知质量；对语义重建而言，遮挡、投影错位以及标签中对索间区域的二维定义又会造成不同的误差结构。若分别训练 RGB 场和二维分割器，输出难以保持跨视角几何一致性；若只把二维预测直接写入高斯属性，又容易把编码器误差固化到错误的三维位置。

本文的目标不是提出一个脱离已有工作的全新渲染原语，而是把桥梁数据的相机约束、显式语义高斯和 DINOv3 视觉编码器部署路径组织成一条可检查的相机输入到 RGB/ID-mask 输出链路。语义场内部共享三维位置和透明度排序，同时保存 RGB 外观与低维语义特征；最终交付则把该语义场与两个独立 RGB 场组合。投影可靠性、编码器预测熵和多视图软融合属于已实现的研究框架，但并未全部进入最终语义场训练；二维精修头只对三维先验作有界、视图条件的残差。归一化绝对梯度、结构预算和相机残差诊断用于研究和工程筛选，不能直接解释为物理相机校正。

本文的主要贡献和工程交付边界如下：
\begin{enumerate}
  \item 给出一套针对官方桥梁传感器的语义场内部共享几何、最终多场组合的 RGB--语义 Gaussian 管线，包括数据准备、视觉编码器前向、语义解码器训练、渲染和独立 CPU 评分。
  \item 实现 DINOv3 ViT-H+/16 的多层解码器和部署阶段软概率融合，并保留 raw-3D 与最终 ID-mask 两类输出。
  \item 建立 50/41 留出协议、共同原图网格、配对 bootstrap 区间、逐像素混淆矩阵和来源/模型哈希审计，报告正面和负面实验。
  \item 发布最终三场相机条件联合系统、ModelScope 资产快照和 uv 环境；明确外部进度目录的比较限制以及尚未被严格消融支持的机制假设。
\end{enumerate}

\section{相关工作与研究定位}
\subsection{Gaussian 场景表示}
\GS{} 用带有均值、协方差、透明度和视角相关颜色的高斯集合替代显式网格或隐式神经场，在训练和渲染效率之间取得了实用平衡\cite{kerbl2023gaussian}。本文沿用其可微投影和 alpha 合成接口；抗混叠、绝对梯度和 MCMC 式策略均视为已有工程组件，而非独立贡献。与单一 RGB 场相比，本文的语义高斯为每个点增加一个可学习特征向量，经共享分类器得到五类概率，再沿 RGB 的排序和透明度合成。

\subsection{预训练视觉编码器与三维语义}
自监督视觉基础模型能够提供比颜色更稳定的区域表征。DINOv2 证明了无需人工标签也能获得可迁移的密集特征\cite{ocal2023dinov2}；本项目以 ModelScope 中的 DINOv3 ViT-H+/16 权重替代初步方案中的旧骨干，并固定权重 SHA-256。DINOv3 在本系统中是冻结的预训练视觉编码器：它提取多层密集特征，交由可训练语义解码器生成类别概率，并在最终部署阶段对语义场渲染 RGB 进行前向推理；它不读取真实目标照片。Feature 3DGS 已研究将二维基础模型特征写入可渲染高斯场的路线\cite{jiang2024feature3dgs}，本文的区别在于桥梁专用的相机条件推理、官方 ID-mask 输出和多场部署约束。

\subsection{研究边界}
投影不确定性和位姿鲁棒 Gaussian 已有研究基础\cite{zhao2024robust}; 因此本文不宣称首次解决位姿后验、真实标定恢复或通用相机鲁棒性。本文的研究变量是：在已固定官方传感器和有限训练轨迹下，可靠性是否能够改善编码器预测在跨视角之间的传输，以及归一化层选择是否能够改善 RGB 场成员的有效对应。实验结果对前者尚未给出独立、稳定的机制收益，对后者只给出有限的归一化条件对照证据。

\section{问题定义与数据协议}
给定训练相机集合 $\mathcal{C}_{\mathrm{tr}}=\{(K_i,T_i,I_i)\}$ 和有限的语义标注 $Y_i$，目标是学习一个场景状态 $\Theta$，对未提供目标照片的相机 $(K,T)$ 输出
\begin{equation}
  (\widehat I(K,T),\widehat M(K,T))=\mathcal{R}_{\Theta}(K,T),
\end{equation}
其中 $\widehat I\in[0,1]^{H\times W\times3}$，$\widehat M\in\{0,1,2,3,4\}^{H\times W}$。五类 ID 的语义约定为 0 背景、1 桥面、2 拉索、3 桥塔、4 基础。推理阶段只读取模型、训练来源照片/缓存和输入相机，不读取目标 RGB、目标 mask 或历史预测 PNG。

正式采用协议固定官方传感器为 $1320\times989$，内参
\[
K=\begin{bmatrix}925.701619&0&660\\0&925.701619&494.5\\0&0&1\end{bmatrix},\qquad
\mathbf d=(0.008987863,0,0,0,0),
\]
并使用原始畸变图像网格计算最终指标。50 张视角用于 RGB，41 张拥有官方 JSON 标签的视角用于语义；训练使用 350 张 TRAIN 来源照片。该划分与外部进度目录中 RGB148/SEM380 的全量拟合数字不同，不能直接相减并解释为泛化提升。

\section{方法}
\subsection{语义场内部的共享几何与双头属性}
语义场由 $N$ 个高斯组成。第 $j$ 个高斯参数为
\begin{equation}
  g_j=(\boldsymbol\mu_j,\boldsymbol{\Sigma}_j,\alpha_j,\mathbf c_j,\mathbf z_j),
\end{equation}
其中 $\boldsymbol\mu_j\in\R^3$、$\boldsymbol{\Sigma}_j$ 为三维协方差、$\alpha_j$ 为不透明度、$\mathbf c_j$ 为 RGB/球谐外观，$\mathbf z_j\in\R^{16}$ 为语义特征。对于相机射线按深度排序后的高斯序列，颜色渲染为
\begin{equation}
  \widehat{\mathbf C}(\mathbf p)=\sum_{j\in\mathcal{V}(\mathbf p)}w_j(\mathbf p)\,\mathbf c_j(\mathbf v),
  \quad w_j=\alpha_j(\mathbf p)\prod_{k<j}(1-\alpha_k(\mathbf p)).
\end{equation}
语义场先通过共享分类器得到 $\mathbf q_j=\operatorname{softmax}(W\mathbf z_j+b)$，再用同一组 $w_j$ 合成 $\mathbf p_{3D}(\mathbf p)=\sum_jw_j\mathbf q_j$。因此，在这个语义场内部，RGB 与 raw-3D 语义共享可见性；最终部署的两个独立 RGB 场不与该语义场共享可训练高斯参数。

\subsection{DINOv3 视觉编码器、语义解码器与训练边界}
冻结 ModelScope DINOv3 ViT-H+/16，参数规模约 $8.41\times10^8$。从多个 transformer 层抽取特征，经线性投影、上采样和 DPT 风格融合得到语义 logits；同时保留浅层 RGB 细节支路和边界头。记经视觉编码器和语义解码器计算得到的类别概率为 $p_{\mathrm{enc}}=\operatorname{softmax}(h_{\mathrm{enc}}(I))$。DINOv3 主干保持冻结，训练参数只属于语义解码器和后续精修头，系统未设置额外语义网络或特征/概率对齐损失。最终深度矩交叉增强语义场的实际训练配置必须单独说明：\texttt{pseudo\_dir: null}、\texttt{pseudo\_weight=0}、\texttt{multiview\_fusion=false}、\texttt{multiview\_weight=0}、\texttt{semantic\_geometry\_weight=0}，并冻结几何和 RGB。因此，多视图预测监督和语义几何联合损失不是该最终语义场高分的直接来源。

部署阶段的路径不同于训练阶段。对输入相机，先渲染深度矩交叉增强语义场自身的 RGB 和 $p_{3D}$；再把该 RGB（不是两个独立 RGB 场的最终融合图）送入 DINOv3 H+ 及语义解码器。推理使用 768 像素滑窗、512 stride、水平翻转 TTA，以及 0.75 局部滑窗概率与 0.25 全景上下文概率的固定融合，最后将编码器—解码器概率与语义场概率按冻结的 0.5/0.5 规则对齐回原始畸变网格并取 ID。测试相机不需要真实目标照片，但需要 DINOv3 H+ 权重、语义解码器和对应运行环境。

\subsection{研究框架中的投影可靠性与多视图融合}
相机和点的局部误差通过投影 Jacobian 传播为像素椭圆的质量代理。对编码器—解码器概率计算熵 $H(p_{\mathrm{enc}})$、可见性、跨视角一致性和投影覆盖度，形成软权重
\begin{equation}
  \rho_i(\mathbf p)=\exp(-\tau H(p_{\mathrm{enc},i}))\,v_i(\mathbf p)\,s_i(\mathbf p),
  \qquad
  \widetilde p_{\mathrm{enc}}=\frac{\sum_i\rho_i p_{\mathrm{enc},i}}{\sum_i\rho_i+\varepsilon}.
\end{equation}
$\rho$ 用于研究框架中的样本接受与监督强度，并不被解释为经过覆盖率校准的完整 BA 后验。下式表示已实现的通用候选目标；最终深度矩交叉增强语义场采用上段列出的 GT-only/冻结场配置，使编码器概率监督与多视图项的系数均为零：
\begin{equation}
\mathcal L_{sem}=\mathcal L_{CE}(p_{2D},Y)+\lambda_T\,\mathrm{KL}(\widetilde p_{\mathrm{enc}}\|p_{2D})
+\lambda_3\,\mathcal L_{CE}(p_{3D},Y)+\lambda_b\mathcal L_{boundary}.
\end{equation}

\subsection{有界视图条件精修}
最终输出还使用一个输入为 RGB、深度、alpha、渲染特征和 $p_{3D}$ 的轻量多尺度头。为限制二维头完全覆盖三维先验，输出写为
\begin{equation}
  p_{final}=\operatorname{softmax}(\log(p_{3D}+\epsilon)+r),
  \qquad |r_c|\le B,
\end{equation}
其中默认 $B=6$。该设计使最终掩码可以表达视图条件的边界和遮挡，同时保留可解释的显式三维概率。报告同时保存 $p_{3D}$ 与 $p_{final}$；后者的高分不能被描述为 raw 三维语义精度。

\subsection{IBGS 来源层条件 RGB 精修}
最终多场系统的 IBGS 分支不是将一个邻居 RGB 当作单个 Gaussian 的颜色，而是把多个来源视图中的层证据传给残差网络。对每个目标相机，在严格距离 $0.01<d<1.5$、方向夹角小于 $30^\circ$ 的 TRAIN 候选中选择至多四个邻居；350 张 TRAIN 缓存中 344 张拥有四个邻居，6 张不足四个，不复制来源填槽。每个目标像素固定保留四个目标层、四个来源视图，即四层 $\times$ 四来源的结构。

对层 $i$ 和来源 $s$，七维输入为
\begin{equation}
 x_{is}=[\Delta RGB_{is},\Delta t_{is},\cos(\theta_{is})]\in\R^7,
\end{equation}
其中 $\Delta RGB$ 是来源 RGB 与目标基础渲染 RGB 的差值，$\Delta t$ 是相机中心位移，$\cos(\theta)$ 是目标/来源射线夹角余弦。来源的原始贡献质量是选择器提供的 $w_i=\alpha_i T_i$，不先归一化；同 ID 的来源支持 $q_{is}\in[0,1]$ 单独记录。逐层输入先经过复用的 $7\to32\to32$ MLP，再按 $w_iq_{is}$ 汇聚四层四来源特征：
\begin{equation}
 z(u)=\frac{1}{4}\sum_{i=1}^{4}\sum_{s=1}^{4}w_i(u)q_{is}(u)\phi(x_{is}(u)).
\end{equation}
归一化对照把分母改为有效支持总质量；无正支持像素在 MLP 前置零、MLP 后再次乘支持，最终残差强制为零。汇聚的 32 通道与目标射线 RGB 和基础 RGB 拼接为 38 通道，再送入空间 CNN 输出三通道颜色残差。由于来源 RGB 是多层 alpha 混合观测，$q$ 只表示同 ID 的贡献支持，不等于物理表面颜色、遮挡真值或单个 Gaussian 的纯辐射。

\subsection{深度矩与交叉特征}
为给语义精修头提供可解释的遮挡深度统计，所有矩使用同一 alpha 合成权重 $a_j=\alpha_j\prod_{k<j}(1-\alpha_k)$。对深度 $z_j$ 和 16 维语义特征 $\mathbf f_j$，在同一可见性序列上计算
\begin{equation}
 E[z]=\frac{\sum_j a_jz_j}{\sum_ja_j},\quad E[z^2]=\frac{\sum_ja_jz_j^2}{\sum_ja_j},\quad E[\mathbf f]=\frac{\sum_ja_j\mathbf f_j}{\sum_ja_j},\quad E[z\mathbf f]=\frac{\sum_ja_jz_j\mathbf f_j}{\sum_ja_j}.
\end{equation}
因此方差通道为 $E[z^2]-E[z]^2$，交叉通道为 $E[z\mathbf f]-E[z]E[\mathbf f]$；深度先按固定场景尺度归一化，交叉项再作有界的 asinh 压缩。新增通道总数为 17：1 个深度方差加 16 个深度--特征交叉量。固定相同 warm-start 和 3,000 步精修预算比较 zero（17 个零输入）、variance（仅方差）和 cross（方差加 16 个交叉量）三臂。cross 相对 variance 的全类 mIoU 差为 $-0.0384$ 个百分点，95\% 区间 $[-0.2049,+0.1796]$，没有稳定收益；因此不能把交叉项写成已验证创新。

方差消减在远距离单层时会受到 FP32 舍入影响。实现使用 $8\epsilon\max(|E[z^2]|+E[z]^2,1)$ 的 floor，并要求 $\alpha\ge10^{-4}$ 且方差超过 $\max(\mathrm{floor},10^{-6})$；这是数值退化保护，不是新的统计校准。

\subsection{归一化梯度与结构预算}
RGB 和语义的高斯更新采用按参数块归一化的绝对梯度统计，减少高分辨率、颜色尺度和语义通道之间的数量级偏差。每次增密将候选分为保留、克隆和结构分裂，并受最大高斯数与最小多视图支持约束。相机残差模块只在训练诊断中计算六维小扰动的有限差分 Jacobian，求带先验的受限最小二乘候选；验证重渲染失败时丢弃候选。测试相机不做在线优化。

\subsection{最终三场组合}
最终三场相机条件联合系统使用三套独立场：约 996,009 点的 IBGS RGB 场、498,136 点的深度矩交叉增强语义场和 500,000 点的 MCMC RGB 场。语义场内部共享其 RGB 与语义几何；最终 RGB 则由两个独立 RGB 场按冻结成员组合产生。DINOv3 语义前向始终读取语义场自身渲染 RGB，不读取 RGB 融合结果。该多场组合是工程部署策略，不被写成新的神经渲染原语。

\begin{figure}[t]
\centering
\begin{tikzpicture}[scale=.72,transform shape,node distance=7mm and 9mm,>=Latex,font=\small]
\tikzset{box/.style={draw,rounded corners,align=center,minimum height=9mm,minimum width=27mm,fill=blue!5}, arrow/.style={->,thick}}
\node[box] (data) {官方相机\\+ TRAIN RGB/标签};
\node[box,right=of data] (geo) {SfM/几何初始化\\共享高斯位置};
\node[box,above right=of geo] (encoder) {冻结 DINOv3 H+\\多层特征+解码器};
\node[box,below right=of geo] (rgb) {RGB 场\\SH/颜色头};
\node[box,right=29mm of geo] (fusion) {投影可靠性\\多视图软融合};
\node[box,right=of fusion] (joint) {语义高斯\\raw-3D + 有界精修};
\node[box,right=of joint] (render) {新相机\\RGB / ID mask};
\draw[arrow] (data)--(geo);
\draw[arrow] (geo)--(encoder);
\draw[arrow] (geo)--(rgb);
\draw[arrow] (encoder)--(fusion);
\draw[arrow] (rgb)--(fusion);
\draw[arrow] (fusion)--(joint);
\draw[arrow] (rgb.east)--(joint.west);
\draw[arrow] (joint)--(render);
\end{tikzpicture}
\caption{训练与推理数据流：语义场内部共享 RGB 与语义几何，最终 RGB 来自两个独立 RGB 场；DINOv3 编码器支路表示部署阶段前向。测试相机不需要真实目标照片，但仍需 DINOv3 H+ 及解码器。}
\label{fig:pipeline}
\end{figure}

\section{实现与复现}
\subsection{软件与资产}
代码使用 Python 3.11--3.12，依赖通过 uv 和提交的 \texttt{uv.lock} 管理，核心组件包括 PyTorch 2.8、gsplat 1.5.3、Transformers、ModelScope、NumPy 和 LPIPS。数据集、DINOv3 权重和训练检查点不放入 Git；它们发布在 ModelScope 数据集\href{https://modelscope.cn/datasets/sky931/SHM2026/files}{sky931/SHM2026}，仓库提供下载和 SHA-256 校验脚本。该组织方式避免 Git 历史被多 GB 二进制污染，也使数据、DINOv3 编码器、语义解码器和最终联合系统组件能够按清单复现。

\subsection{训练阶段}
训练首先建立 RGB/几何场，再从 warm-start 检查点进入语义阶段。最终深度矩交叉增强语义场的 8,000 步配置冻结 geometry 与 RGB，使用 259 张标注 TRAIN 图和多尺度有界精修头；\texttt{pseudo\_dir: null}、\texttt{pseudo\_weight=0}、\texttt{multiview\_fusion=false}、\texttt{multiview\_weight=0} 和 \texttt{semantic\_geometry\_weight=0}。随后固定场的 3,000 步深度矩对照只更新精修头，比较 zero、variance、cross 三臂。DINOv3 编码器前向发生在最终部署；最终语义场训练不把编码器概率作为二维预测监督的来源。每个阶段保存配置、输入哈希、checkpoint 和评价回执；正式验证图只在冻结预测落盘后被评分进程读取。

表~\ref{tab:active-config} 将研究框架开关与最终语义场实际值分开列出。
\begin{table}[t]
\centering
\caption{最终深度矩交叉增强语义场的关键配置。}
\label{tab:active-config}
\begin{tabular}{lll}
\toprule
配置项 & 最终值 & 解释\\
\midrule
\texttt{pseudo\_dir} & \texttt{null} & 不读取预存二维预测\\
\texttt{pseudo\_weight} & 0 & 不加入额外二维预测监督\\
\texttt{multiview\_fusion} & \texttt{false} & 不做训练期多视图软融合\\
\texttt{multiview\_weight} & 0 & 多视图项无训练权重\\
\texttt{semantic\_geometry\_weight} & 0 & 语义阶段不更新几何\\
\texttt{freeze\_geometry/rgb} & \texttt{true/true} & 仅更新语义/精修参数\\
\bottomrule
\end{tabular}
\end{table}

\begin{algorithm}[t]
\caption{相机输入到双输出的部署流程}
\label{alg:deploy}
\begin{algorithmic}[1]
\Require 输入相机 $(K,T)$，三场部署清单，TRAIN 来源缓存，DINOv3 H+ 与语义解码器
\Ensure RGB PNG 和单通道 ID-mask PNG
\State 校验 bundle、模型和来源缓存 SHA-256
\State 依据距离和方向阈值选择至多四个 TRAIN 邻居
\State 对 IBGS、语义场、MCMC 三场执行可见性排序和 alpha 合成
\State 融合两个独立 RGB 场得到最终 RGB，同时保留语义场自身 RGB
\State 将语义场自身 RGB 以 768/512 滑窗、翻转 TTA 和局部/全景融合送入 DINOv3 H+
\State 将编码器—解码器概率与语义场 $p_{3D}$ 按冻结 0.5/0.5 规则计算 $p_{final}$
\State 将 $\arg\max_c p_{final,c}$ 映射为官方 0--4 ID，保存 PNG
\State 禁止读取目标照片、目标标签和历史预测；不优化测试相机
\end{algorithmic}
\end{algorithm}

\section{实验设置}
\subsection{评价指标}
RGB 使用原图像素的
\begin{equation}
\PSNR=10\log_{10}\frac{1}{\mathrm{MSE}},
\end{equation}
并报告 SSIM 和 LPIPS（后者越低越好）。语义报告五类 pooled confusion matrix 的平均 IoU：
\begin{equation}
\mIoU=\frac{1}{5}\sum_{c=0}^{4}\frac{TP_c}{TP_c+FP_c+FN_c}.
\end{equation}
此外保存前景 mIoU、拉索 IoU、边界 F1 和 raw-3D mIoU。所有正式差异按同名相机配对，并使用 5000 次 bootstrap 的双侧 95\% 区间；独立 CPU 审计以 LabelMe 顺序重绘标签、重建逐图混淆矩阵和 RGB 指标。

\subsection{主结果}
\begin{table}[t]
\centering
\caption{共同原图协议上的正式留出结果。RGB 为 50 视角，语义为 41 视角。}
\label{tab:main}
\begin{tabular}{lrrrr}
\toprule
系统 & PSNR $\uparrow$ & SSIM $\uparrow$ & LPIPS $\downarrow$ & 五类 $\mIoU\uparrow$\\
\midrule
本地 SEM386-style Swin-L 参考 & 29.400950 & 0.870650 & 0.271237 & 94.318382\%\\
基准三场联合系统 & 30.035661 & 0.874421 & 0.264742 & 95.109016\%\\
\textbf{最终三场相机条件联合系统} & \textbf{30.236593} & \textbf{0.875024} & \textbf{0.259300} & \textbf{95.109016\%}\\
\bottomrule
\end{tabular}
\end{table}

最终三场相机条件联合系统相对本地参考的配对差异为 PSNR +0.835644\,dB（95\% 区间 [0.564120,1.139829]）、SSIM +0.004374、LPIPS -0.011937、\mIoU{} +0.790633 个百分点（区间 [0.136189,1.653615]）。相对基准三场联合系统，PSNR +0.200932\,dB（区间 [0.111521,0.307454]），SSIM +0.000603，LPIPS -0.005442，而语义混淆矩阵逐字节不变。这说明 RGB 成员替换在当前协议下没有损伤已选语义，但不能把固定组合的收益归因于某一个单场机制。

\subsection{与外部进度目录的协议分层对照}
\begin{table}[t]
\centering
\caption{与 \texttt{project\_progress\_20260918} 的协议分层对照。}
\label{tab:peer}
\begin{tabular}{lrrr}
\toprule
条目与协议 & RGB PSNR & RGB SSIM & 五类 $\mIoU$\\
\midrule
外部目录 RGB148+SEM380，全量拟合 300 视角 & 23.305737 & 0.896586 & 95.306455\%\\
最终联合系统，全量复核 300 个有标签视角 & \textbf{31.095728} & 0.876151 & \textbf{96.590043\%}\\
外部目录 SEM386 Dev30 留出（记录值） & -- & -- & 93.7418\%\\
最终联合系统，正式 50/41 留出 & 30.236593 & 0.875024 & 95.109016\%\\
\bottomrule
\end{tabular}
\end{table}

全量表显示最终联合系统在 PSNR 和 mIoU 上高于外部进度目录记录，但 SSIM 较低，而且训练视角、检查点和预算并不完全匹配，因此它只是工程复核对照。正式留出结果优先于全量拟合数字；对外部 Dev30 只能说明数值上高于公开记录，不能宣称精确复现或赛事盲测胜出。

\subsection{消融与负面结果}
\begin{table}[t]
\centering
\caption{关键候选的终点评价与采用决定。未采用结果保留用于防止选择性报告。}
\label{tab:ablations}
\scriptsize
\begin{tabularx}{\linewidth}{>{\raggedright\arraybackslash}p{3.3cm} X r r l}
\toprule
候选 & 唯一主要变化 & PSNR & 五类 $\mIoU$ & 决定\\
\midrule
语义场与 DINOv3 固定部署组合 & 语义场概率与编码器—解码器概率各半 & 29.400950 & 95.109016\% & 基准联合系统\\
归一化 top4 与匹配置换控制 & 正确对应相对训练置换控制 & +0.045201 dB（相对控制） & -- & 有限机制证据\\
渲染 RGB 域解码器适配 & 350 TRAIN 渲染图继续训练 2000 步 & -- & 94.415452\% & 低于最终系统\\
固定 RGB 加权 $\alpha=.6$ & 调整两 RGB 场权重 & 30.252777 & -- & PSNR 门失败\\
warm-IBGS 端口 & 来源输入与稳定 AA 训练 & 29.846757 & -- & 低于基准系统\\
DINOv3 7B 容量对照 & 更大视觉编码器容量 & -- & 95.18771\% & 区间跨零\\
连续语义几何梯度 & joint/PCGrad & -- & 94.932812/94.930343\% & 最终退化\\
\bottomrule
\end{tabularx}
\end{table}

正确层对应相对匹配训练的特征置换控制提升 0.045201 dB，95\% 区间 [0.020287,0.068611]，SSIM 提升 0.000828；LPIPS 区间跨零。这个结果与四组层选择主对照不同：在同一归一化汇聚条件下，top4-normalized 相对 median4-normalized 的 PSNR 提升 0.127846 dB、SSIM 提升 0.001250559、LPIPS 降低 0.003953302，三个配对区间均指向改善。前者是“正确对应 vs 置换控制”，后者是“top4 vs median4”的层选择对照；两者不能互换，也不能包装为完整系统的独立主贡献。渲染 RGB 域适配、7B 编码器扩容、普通域适配、连续语义位置修正和多个结构交换均未通过预设采用门；这些负结果是报告研究边界的重要组成部分。

\subsection{成本、鲁棒性与失败模式}
最终三场相机条件联合系统包含三套独立场，共 1,994,145 个高斯：IBGS 996,009、深度矩交叉增强语义场 498,136、MCMC 500,000；另有约 66,095 个 RGB 头参数、840,592,640 参数的冻结 DINOv3 H+ 视觉编码器、350 张来源照片及约 21.93 GB 层缓存。52 个相机（50 正式视角，加改名和位姿压力视角）共完成 156 次场渲染；完整导出约 118.2 秒，峰值 CUDA 分配约 7.81 GB。推理只支持训练银行对应的官方传感器和至少一个满足距离/方向阈值的 TRAIN 邻居；没有邻居时明确失败，不承诺任意相机内参或场外位姿。

实际核验包括：改名不影响来源选择；对同一相机施加 $w2c[0,3]+0.1$ 后 RGB 和 mask 均改变；mask 只含 0--4；50 张 RGB 和 50 张 mask 重新落盘并经独立 CPU 复核。部署输入必须为固定官方传感器尺寸 $1320\times989$、对应内参和畸变参数，只允许改变位姿；至少需要一个满足上述距离/方向门限的 TRAIN 邻居。测试阶段不优化相机、不读取目标 RGB，也不使用历史预测 PNG。只复制一个 checkpoint 无法部署：必须同时提供三个场、DINOv3 H+ 编码器权重和语义解码器、350 张来源照片、约 21.93 GB 层缓存、冻结 CUDA/Python 源码以及两个 uv 运行环境。早期冻结编码器接口缺失导致的 v1 失败保留在运行回执中，v2 仅修正多余接口调用后重新执行，未调整模型数学。

\section{讨论}
\subsection{结果含义}
当前结果表明，在严格固定的官方传感器和共同原图网格上，多场 RGB 融合与语义场+DINOv3 部署路径能同时提供较高 RGB 质量和五类掩码质量。共享可见性只适用于语义场内部；最终联合系统的 RGB 来自两个独立 RGB 场。有界精修允许最终 mask 处理视图条件边界，而 raw-3D 输出保持对显式三维属性的审计。对工程部署而言，ModelScope 清单、模型哈希和 uv 锁定环境比单个 checkpoint 更重要，因为最终系统依赖来源银行、缓存、三个场和两个运行环境的协同。

\subsection{学术创新的诚实边界}
本文最有研究价值的假设是“相机/投影可靠性可以指导编码器—解码器预测在跨视角传输中的权重”和“归一化特征层对应可影响多场 RGB 成员质量”。现有实验只为后者提供了小幅、受条件限制的正证据；前者的多数长训练候选没有超过最终联合系统。DINOv3、DPT 解码器、EMA、一致性损失、Gaussian rasterization、绝对梯度和固定成员平均都已有先例或属于工程组合。因而本报告将最终三场相机条件联合系统定义为高质量、可复现的工程方案，并把机制创新表述为待进一步检验的研究假设，而非“首次”结论。

\subsection{局限与未来工作}
首先，实验集中于单桥和单一官方传感器，不能推断跨桥梁、跨焦距或跨天气泛化。其次，三场与 21.93 GB 缓存带来较高存储和部署成本；只复制一个 checkpoint 不能运行，还必须迁移来源照片、层缓存、DINOv3 H+、语义解码器及两个 uv 运行环境。第三，最终 mIoU 依赖二维精修，raw-3D mIoU 明显更低，说明显式三维语义仍有改进空间。第四，外部进度目录缺少完整 Dev30/RGB150 检查点和清单，严格公平重训尚未完成。后续应在多个桥梁、匹配训练预算和固定随机种子下，补齐投影各向同性对照、同一编码器不同层对应、三场成员的单变量训练以及 raw/final 双重语义指标。

\section{结论}
本文完成了一个从官方相机输入到 RGB 与语义 ID-mask 输出的 Bridge-RGS 研究型实现，并以 ModelScope 管理数据、DINOv3 权重和训练检查点。最终三场相机条件联合系统在 50/41 共同原图留出协议上达到 30.236593 dB PSNR、0.875024 SSIM、0.259300 LPIPS 和 95.109016\% 五类 mIoU；相对本地同协议参考，RGB 与语义均有明确提升。报告同时保留了全量拟合、负面消融、资源成本和协议限制。该结果足以作为公开仓库和比赛工程交付；若要形成可发表的学术主张，还需要跨场景、匹配预算和更严格的机制隔离实验。

\section*{可复现清单}
\begin{itemize}
  \item 代码仓库：\href{https://github.com/QinnyuLan/SHM26P2}{github.com/QinnyuLan/SHM26P2}。
  \item 大型资产：\href{https://modelscope.cn/datasets/sky931/SHM2026/files}{ModelScope sky931/SHM2026}，含 \texttt{dataset/}、\texttt{models/dinov3-vith16plus/}、\texttt{checkpoints/} 和 \texttt{manifest.sha256}。
  \item 下载命令：\par\texttt{uv run python scripts/download\_modelscope\_assets.py}\par\texttt{--output release\_assets/SHM2026}。
  \item 部署入口：\par\texttt{uv run bridge-rgs render --checkpoint ...}\par\texttt{--cameras ... --output ...}；三场完整入口和输入限制见\par\texttt{docs/ibgs\_joint\_deployment.md}。
  \item 主结果与协议：\par\texttt{docs/ibgs\_joint\_bundle\_results.md}\par\texttt{docs/project\_progress\_comparison.md}\par\texttt{docs/experiment\_history\_20260927.md}。
\end{itemize}

\bibliographystyle{plain}
\bibliography{refs}
\appendix
\section{符号与实现对应}
\begin{table}[h]
\centering
\begin{tabular}{ll}
\toprule
符号 & 工程含义\\
\midrule
$\boldsymbol\mu,\boldsymbol{\Sigma},\alpha$ & Gaussian 的位置、协方差和不透明度\\
$\mathbf c$ & RGB/球谐颜色属性\\
$\mathbf z$ & 16 维语义特征\\
$p_{3D}$ & 共享可见性合成的 raw 三维语义概率\\
$r$ & 有界二维视图条件精修残差\\
$\rho$ & 编码器预测熵、可见性和投影支持构成的软可靠性权重\\
\bottomrule
\end{tabular}
\end{table}

\section{报告中的指标解释约定}
所有带有“正式联合系统结果”的数值都来自冻结预测落盘后再读取 GT 的流程；“全量复核”仅用于解释外部进度目录中的展示数字，不能替代留出。不同训练历史、不同视角集合或不同渲染网格之间不得直接作因果差分。负面结果保留是为了防止把单个开发集最高点误写成稳定泛化收益。
\end{document}
